%!TEX root = ../../main.tex
\section{결과의 연결}
\label{sec:disc-chain}

네 연구 질문의 답은 한 사슬을 이룬다. 교전 단위는 자료에서 추정한 $G$, $D$와 규칙에 고정한 $R$, $B$로 만들어지고, 작은 쪽 참여 인원으로 한타 T와 소규모 교전 S로 나뉜다(M-RQ1). 이 단위가 입력(예측 시점 $\tcut$의 상태)과 결과(결과 시점 $\eh$의 상태)를 함께 정하며, 정의 상수를 일변량으로 바꾸면 존재 게이트의 두 상수가 사례 집합을 가장 크게 움직인다. 결과 가치는 그 두 시점에 동결 평가기를 적용한 변화의 부호이고(M-RQ2), 평가기 품질, 삼중 분해, 비교전 대조, 물질 대응, 지평·동료 평가기 안정성, 시간 경과와 관측 갱신의 분해가 그 측정이 무엇을 재는지 한정한다. 사전 예측은 코호트마다 하나의 고정 사양 예측기를 승률·시간 스플라인 기준선과 같은 행에서 대비하여 초기 우세와 시간을 넘는 이득을 보였고, 정보군 비교는 사건 특징을 프레임 특징에 추가하는 효용을 두 코호트에서 보였다(M-RQ3). 검증은 그 이득이 균형 구간에서 더 작고, 작은 변화 제외에 견디며, 판별 성분과 함께 나타나고, 새 프레임이 들어온 셀에서 전체와 같은 부호이며, 평가한 외부 표본에서는 한타의 우세가 관측되지 않았음을 보였다(M-RQ4).

각 단계의 해석 조건은 다음 단계로 상속된다. 정의가 사람의 주석으로 검증되지 않았으므로 $\SVI$는 검출된 교전에 대한 것이고, $\SVI$가 모델 정의이므로 $q$의 이득은 그 평가기에 조건부이며, 외부 평가가 점수 전용이므로 외부에서 달라진 성능의 원인은 식별되지 않는다. M-RQ3의 두 번째 절은 둘로 나뉘어 답해졌다. 관측 정보의 구성은 동일 학습 절차의 정보군 비교로 답했고, 학습기 선택의 독립적 영향은 동일 입력의 정합 비교가 미실행이므로 고정 사양 아래의 결과로만 답했다.
